\chapter{共轭梯度法}\label{v0ch:13}
\begin{chaptersummary}
一般非对称问题需要长基或重启，而 SPD 结构允许以短递推维持共轭方向。本章由二次能量最小化组织 CG，明确其正定性条件与误差界，并给出完整两步算例。后面的 SPD 预条件和多重网格组合必须保持这些结构。
\end{chaptersummary}
\learninggoals{执行 CG 与 PCG 递推；解释 A-共轭和能量误差；识别标准 PCG 条件。}
\paragraph{核心问题}SPD 结构怎样减少存储，又保留子空间最优性？

\section{能量最小化与短递推}
对 SPD $A$，CG 在 $\bm x_0+\K_k(A,\bm r_0)$ 中最小化
$\norm{\bm x_*-\bm x}_A$，其中 $\norm{\bm e}_A^2=\bm e^TA\bm e$。
初始化 $\bm p_0=\bm r_0$，递推
\[
 \alpha_k=\frac{\bm r_k^T\bm r_k}{\bm p_k^TA\bm p_k},\quad
 \bm x_{k+1}=\bm x_k+\alpha_k\bm p_k,\quad
 \bm r_{k+1}=\bm r_k-\alpha_kA\bm p_k,
\]
\[
 \beta_k=\frac{\bm r_{k+1}^T\bm r_{k+1}}{\bm r_k^T\bm r_k},\qquad
 \bm p_{k+1}=\bm r_{k+1}+\beta_k\bm p_k.
\]
精确算术下，残差彼此正交，方向 $A$-共轭；正定性保证每个非零搜索方向的分母为正，支撑这一短递推。
\section{两步算例与误差界}
取 $A=\left[\begin{smallmatrix}4&1\\1&3\end{smallmatrix}\right]$、
$\bm b=(1,2)^T$、$\bm x_0=0$。
第一步 $\alpha_0=1/4$，
$\bm x_1=(1/4,1/2)^T$，
$\bm r_1=(-1/2,1/4)^T$；
$\beta_0=1/16$，
$\bm p_1=(-7/16,3/8)^T$。
第二步 $\alpha_1=4/11$，得到
$\bm x_2=(1/11,7/11)^T$。
经典界为
\[
 \frac{\norm{\bm e_k}_A}{\norm{\bm e_0}_A}
 \le2\left(\frac{\sqrt\kappa-1}{\sqrt\kappa+1}\right)^k,
 \qquad \kappa=\kappa_2(A).
\]
这是最坏情形上界；特征值聚簇与右端方向可使实际收敛更快，二范数残差也不保证逐步单调。
\section{PCG 的正确条件}
取固定 SPD $M$，解 $\bm z_k=M^{-1}\bm r_k$。将 $\alpha$ 分子改为 $\bm r_k^T\bm z_k$，$\beta$ 改为相邻该内积之比，方向更新为 $\bm z_{k+1}+\beta_k\bm p_k$。有效 SPD 算子是 $M^{-1/2}AM^{-1/2}$，并非要求 $M^{-1}A$ 在 Euclidean 内积下对称。
\conceptbridge{原矩阵 SPD，为什么任意预条件后不能直接运行标准 CG？}{预条件改变了算法使用的几何。固定 SPD 逆作用对应一个合法能量坐标变换；非对称或随迭代改变的作用通常破坏短递推依赖的正交关系。}{原能量 $\longrightarrow$ 对称正定变换 $\longrightarrow$ 可保持的共轭结构。}
\costnote{每步一次 $A\bm p$、少量向量更新与点积，存储为 $O(N)$；分布式归约延迟可能主导时间。选用 ILU 或 AMG 时，应按标准 PCG 的固定、对称和正定条件检查其完整逆作用。}
\coderef{sparse_linear_algebra/solvers.py}
\aiexercise{手算两步 CG 的全部内积，验证两方向 A-共轭，并比较能量误差与二范数残差。}{用 IC 与非对称 ILU 分别预条件同一 SPD 系统，检查算法条件，不以一次成功运行代替适用性判断。}
\studyroute{略去有限精度共轭丢失及流水化 CG 稳定性证明。参考 \cite{saad,templates}。}{conjugate gradient; A norm; PCG; SPD preconditioner; finite precision CG}
